\documentclass[10pt]{article}
\usepackage[margin=0.8in]{geometry}
\usepackage{amsmath,amssymb,booktabs,graphicx,hyperref}
\setlength{\emergencystretch}{1em}
\newcommand{\pending}{\textemdash}
\title{Ray-Based Semantic Evidence Retargeting for Multi-View 3D Editing}
\author{Working manuscript; title and final model subject to confirmation}
\date{Development draft, 24 October 2026 UTC}
\begin{document}
\maketitle

\noindent\textbf{Draft status.} Method and protocol text below describe implemented
candidate designs. The final architecture is not yet selected. All populated
results are development measurements, not held-out test results. Empty entries
remain unmeasured. Core related-work references have been checked against
primary sources; the broader literature review remains open. The final
abstract and generalization claims require the planned confirmation experiments.

\section{Problem and design motivation}
We study multi-view 3D editing from one edited target RGB image and three
undeformed source RGB images with calibrated cameras. The target specifies a
geometric edit, while the sources provide complementary appearance evidence.
Unchanged spatial projection can associate a source feature with an incorrect
location after deformation. Semantic retrieval offers an alternative, but a
single retrieved correspondence can be unreliable, and a better correspondence
score does not by itself establish an improvement in the generated asset.

We investigate a source-adaptation path that preserves the pretrained
multi-view conditioning function at initialization. The design separates
candidate retrieval, candidate aggregation, ray-based feature delivery, and
retargeting. This separation permits matched comparisons of individual choices
and prevents a gain from one component from being attributed to another.

\section{Related work}
\paragraph{Image-conditioned 3D generation.}
TRELLIS.2 learns a compact native representation of geometry and material
attributes using sparse voxels~\cite{xiang2025trellis2}.
Pixal3D lifts image features into a back-projected conditioning volume and
extends this mechanism to multiple views~\cite{li2026pixal3d}.
Our implementation continues the supplied multi-view Pixal3D checkpoints
and adapts their source-conditioning path. The separately implemented
multi-view prior and checkpoint lineage are specified in our protocol;
we do not identify every implementation detail with the original paper branch.

\paragraph{Image-conditioned editing.}
ShapeUP learns supervised latent translation from a source 3D shape and an
edited image to an edited 3D shape~\cite{gat2026shapeup}.
Our evaluated input instead consists of an edited target image and undeformed
source images with calibrated cameras. This input distinction motivates our
study of appearance evidence under conflicting geometric states. It does not
establish a performance advantage over source-shape-conditioned editing;
no such quantitative comparison has yet been executed in this study.

\paragraph{Representations for correspondence.}
Zhang et al. examine complementary correspondence properties of diffusion
and DINO features~\cite{zhang2023twofeatures}.
VGGT-$\Omega$ learns geometry-aware representations through reconstruction
training~\cite{wang2026vggtomega}. RoMa-$\Omega$ distinguishes zero-shot
feature matching from learned matching on feed-forward 3D
representations~\cite{nordstrom2026romaomega}.
These works motivate separating raw descriptor quality, learned matching,
and candidate aggregation. Our experiments additionally ask whether their
correspondence evidence improves generated geometry and source appearance
when delivered through the pretrained projection-conditioning path.


\section{Method}
\subsection{Pretrained multi-view carrier and ray groups}
Let $T$ be the native target feature map, $S_j$ the native feature map of source
view $j\in\{1,2,3\}$, and $P_j$ its calibrated projection into the SS conditioning
grid. We retain the original four-view carrier $C_{\mathrm{MV4}}$, including
its native target contribution and global-conditioning tokens. For a source
camera and image patch $p$, define a ray group
\begin{equation}
 \mathcal R_{j,p}=\{x:\pi_j(x)=p\}.
\end{equation}
Voxels assigned to this camera ray share an image-space evidence decision.
Native bilinear projection distributes the resulting map to the grid. This is
different from predicting an independent attention output for each voxel.
The grouping is determined by the input cameras; it does not require a target
surface depth or a surface-visibility test at inference.

\subsection{Semantic proposals and member-preserving matching}
For each observed target patch $m$, raw semantic similarity retrieves an
initial set of $K=64$ source patches across the three source images. A frozen
matching decoder contextualizes the target descriptor, and its score head
produces a categorical distribution $p_m(k)$ on these proposals. We separately
evaluate raw semantic descriptors, learned contextual descriptors, and use of
the scalar matching score; these are distinct design choices.

One source-camera ray may intersect grid sites assigned to different target
patches. Let $\mathcal M_r$ contain the distinct observed target patches of ray
$r$. The member-based candidate construction averages their normalized
distributions before retaining a final $K$-element support:
\begin{equation}
 \bar p_r(k)=\frac{1}{|\mathcal M_r|}\sum_{m\in\mathcal M_r}p_m(k),\qquad
 p_r(k)=\frac{\bar p_r(k)}{\sum_{l\in\mathcal K_r}\bar p_r(l)},\quad k\in\mathcal K_r.
\end{equation}
Here $\mathcal K_r=\operatorname{TopK}(\bar p_r)$; missing entries in an
individual member distribution have zero mass. Duplicate occurrences of the
same target patch vote once. This retains complementary candidate modes
before the final truncation, rather than matching one averaged descriptor.
The observation support uses input image alpha and camera addresses. It is
not an inferred surface-visibility mask.

\subsection{Matching-aware evidence aggregation}
The fixed-evidence candidate aggregates original source values $v_k$ as
\begin{equation}
 E_r=\sum_{k\in\mathcal K_r}p_r(k)v_k.
\end{equation}
The value features retain their native magnitude. The same evidence is used
by the carrier update, source-use decision, and retargeter. Thus matching
weights influence attention without requiring an additional learned value or
output projection.

We also evaluate a learned refinement with four heads and identity values.
For head $h$, the learned score difference is anchored to its initial score:
\begin{equation}
 a_{rhk}=\operatorname{softmax}_{k\in\mathcal K_r}
 [\log p_r(k)+s_{rhk}(\theta)-s_{rhk}(\theta_0)],\qquad
 E_r=\sum_k\left(\frac14\sum_h a_{rhk}\right)v_k.
\end{equation}
The implementation centers the score difference by its prior-weighted mean,
which does not change the softmax. Zero prior support remains zero. This
initializes the learned candidate to exactly the fixed matching distribution.
Additional learned refinement is a tested option, not a prerequisite assumed
by the method description. Comparisons with learned value/output transforms
are reported separately.

\subsection{Source contrast retargeting and native fallback}
Let $u_r$ be a learned source-use probability and $g$ a signed, bounded fusion
gate. The retargeter receives the displacement between retrieved evidence and
the source feature already present in the native carrier, together with
target context:
\begin{equation}
 R_r=R_\phi([E_r-S_r,\,T_r]),\qquad
 \Delta_r=u_r\{g(E_r-S_r)+R_r\}.
\end{equation}
The resulting SS conditioning is
\begin{equation}
 C=C_{\mathrm{MV4}}+\frac14\sum_{j=1}^3P_j(\Delta_j).
\end{equation}
Subtracting $S_r$ is appropriate here because that feature is present in
$C_{\mathrm{MV4}}$. It would not be appropriate for a target-only carrier.
At initialization, $g=0$ and the retargeter output is zero, restoring the
native conditioning function. Rays without matching evidence contribute zero
correction and retain the native multi-view prior. This fallback does not
write a zero-valued observation into the source feature map.

The target feature extraction and its direct conditioning path are retained.
The source branch is the only adapted image-conditioning path. This structural
property does not guarantee exact target RGB preservation by the generated
asset, which remains an empirical evaluation question. No extra
surface-visibility write mask is introduced.

\subsection{Stage composition}
SS64 and PBR1024 continue the supplied multi-view LoRA checkpoints and their
matching configurations. Shape512 and Shape1024 remain frozen. The matching
backbone and decoder are frozen during these generative fine-tuning runs.
PBR retargeting is compared with vanilla PBR on identical generated geometry,
archived native guides, texture sampling noise, and rendering triangles.
The PBR feature-correction candidates are distinct from the SS identity-value
aggregation above; no shared positive effect across stages is presumed.

\section{Experimental protocol}
\subsection{Data and initialization}
Deformation fine-tuning uses the constructed DeformNeXt-10k split, with 9,149
SS training examples, 508 validation examples, and 508 test examples. PBR
training uses 9,139 examples after the recorded training-only exclusion of ten
oversized cases. Source-asset identities are disjoint across these deformation
splits. This does not establish disjointness from the earlier multi-view
pretraining data. The mixed385 subset (80,892 training assets and 950 validation
assets) describes that earlier multi-view fine-tuning, not the deformation
fine-tuning population.

The supplied multi-view checkpoints are SS64 step 4,210 and PBR1024 step
16,850. Both frozen base weights and LoRA weights are loaded using their
original configurations. The frozen shape checkpoints are Shape512 step
4,210 and Shape1024 step 16,850. The matching backbone uses the selected
VGGT-Omega epoch-7 checkpoint. Exact weight and input identities are retained
in the experimental records.

SS comparisons use 480 optimizer updates with global batch 96, LoRA rank and
alpha 16, denoiser learning rate $10^{-5}$, and conditioner learning rate
$10^{-4}$. PBR comparisons use 1,900 updates with global batch 24, denoiser
learning rate $5\times10^{-6}$, and conditioner learning rate $10^{-4}$.
These budgets correspond to approximately five passes through their respective
training populations. Ablation inputs, initialization, budget, sampler, and
evaluation pairing are held fixed. Removed components are physically absent;
absent channels are not replaced by zero placeholders.

\subsection{Development, independent confirmation, and final test}
The SS development panel contains 128 examples from 93 source-asset clusters,
evaluated using three paired generation-noise seeds and 12 sampling steps.
Generation-noise repetitions are averaged within an example before inference
over source clusters. They are not independent training seeds. The completed
SS results in this draft use training seed 101; training seed 202 confirmation
is a separate experiment. Rendering development uses a fixed 32-example
panel and one generation-noise realization.

Final model choice uses both geometric alignment and source preservation,
rather than SS IoU alone. Required validation fresh-load/cache checks,
independent-training confirmation, model selection, and protocol freezing
precede evaluation on all 508 test examples. The final table compares
source-only zero-shot, edited-target zero-shot, vanilla fine-tuning,
retargeting, and the selected combined model when it is distinct. Source-only
zero-shot uses four undeformed M0 RGB views, including the original target
camera. Edited-target input uses M1 RGB at camera 8 and M0 RGB at cameras
9, 10, and 11.

\subsection{Metrics and statistical scope}
Source appearance is evaluated by LPIPS and DreamSim at cameras 9 and 10,
with CLIP as a complementary semantic measure. References preserve the
original material identity on the edited M1 geometry; they are not undeformed
M0 RGB images. Target alignment is evaluated by PSNR, SSIM, and 2D alpha-mask
IoU at camera 8, together with Chamfer distance and surface F-Score in 3D.
Surface evaluation uses 20,000 points and F-Score threshold 0.02 in the fixed
normalized frame, without a result-specific ICP alignment. Voxel IoU is an SS
diagnostic and is not substituted for these mesh/rendering endpoints.

FID and KID assess the global rendered-image distribution using the full
prescribed final panel. They are not per-asset scores. The finite-sample KID
estimate is reported with its sign. Official perceptual/inception evaluation
uses FP32 with TF32 and autocast disabled. All models use the same reference,
camera, rendering, and image preprocessing protocols. Cameras 8--10 are input
camera evaluations; these results are not evidence of novel-view performance.

Paired uncertainty estimates resample source-asset clusters, retaining the
dependence between deformations of an asset and between generation noises.
The original finite ablation family and the subsequent evidence follow-up
are reported as separate studies. The DEV32 intervals in this draft are
exploratory and unadjusted; they are not final test confirmation. Primary
geometry endpoints are CD/F-Score, and primary appearance endpoints are
LPIPS/DreamSim. Positive CLIP results do not change that priority after the fact.

\section{Development results and bounded interpretation}
\begin{table}[t]
\centering\scriptsize
\resizebox{\linewidth}{!}{%
\begin{tabular}{lrrrrrrrr}
\toprule
Actual DEV32 composition&Source LPIPS$\downarrow$&Source DreamSim$\downarrow$&Source CLIP$\uparrow$&Target PSNR$\uparrow$&Target SSIM$\uparrow$&Target M-IoU$\uparrow$&CD$\downarrow$&F-Score$\uparrow$\\
\midrule
Vanilla V1&0.053099&0.155692&0.942434&24.47068&0.950488&0.880945&0.012452&0.849149\\
Top-1 retargeting&0.052104&0.152435&0.940636&24.57466&0.951284&0.885678&0.012413&0.851420\\
D1-cos SS + Top-1 R PBR&0.053362&0.155373&0.942157&24.57729&0.951170&0.886600&0.012504&0.851308\\
Fixed-evidence SS + Top-1 R PBR&0.053290&0.157828&0.941692&24.54433&0.951563&0.886574&0.012328&0.856566\\
Anchored-learned SS + Top-1 R PBR&0.052894&0.156364&0.943132&24.62475&0.951161&0.885769&0.012498&0.854233\\
\bottomrule
\end{tabular}}
\caption{Actual DEV32 compositions, one training seed and one generation-noise realization. The source reference is original material on M1 geometry. The same V1 and retargeting per-asset baselines are verified across the three candidate studies. This is a development table; FID/KID and held-out test results are not included.}
\label{tab:dev-joint32}
\end{table}

\begin{table}[t]
\centering\small
\begin{tabular}{lr}
\toprule
SS model&Macro voxel IoU$\uparrow$\\
\midrule
V1&0.467355\\
C\_SPATIAL&0.471242\\
C\_RAW&0.472360\\
C\_D1\_COS&0.473008\\
C\_DECODER&0.472600\\
C\_DECODER\_NO\_R&0.468671\\
C\_TOP1\_V&0.472218\\
C\_FIXED\_TOPK&0.472745\\
FULL&0.472078\\
FULL\_QK\_ONLY&0.472446\\
FULL\_VOXEL\_Q&0.472314\\
E\_FIXED\_RAW&0.473471\\
E\_LEARNED\_RAW&0.474794\\
FULL\_MEAN\_MATCH&0.472679\\
\bottomrule
\end{tabular}
\caption{Fourteen completed SS endpoints on DEV128, 93 source clusters, three generation noises, training seed 101, and 480 updates. The noises are averaged per example. The original twelve-model comparison family is complete; its ten module contrasts are reported separately. Mean/member is a conditional comparison in the original FULL family. The learned-evidence endpoint belongs to the separately reported exploratory follow-up. FULL is an implementation identifier, not the final selected model.}
\label{tab:dev-ss128}
\end{table}

\begin{table}[t]
\centering\small
\begin{tabular}{lrrrr}
\toprule
Matching diagnostic&Exact@1&Exact@8&Exact@64&Near-1@64\\
\midrule
Raw cosine, mean-ray support&1.25&11.09&50.76&80.41\\
D1 cosine, same support&7.73&28.47&50.76&80.41\\
D1 scalar score, same support&8.56&29.92&50.76&80.41\\
Member-posterior mixture&11.19&45.68&73.36&95.59\\
\bottomrule
\end{tabular}
\caption{Reused matching diagnostics on 50,493 observable source-camera-ray/target-patch relations from 32 development source clusters (percent). Exact denotes the GT source patch on the 32-by-32 image grid; Near-1 allows a one-patch neighborhood. These are conditional relation-coverage measures, not per-point pixel accuracy or generative-quality scores. GT visibility is used only to define the diagnostic denominator, not as an inference write gate.}
\label{tab:matching-diagnostic}
\end{table}


\paragraph{Candidate coverage versus generated quality.}
The reused matching diagnostic contains 50,493 observable
source-camera-ray/target-patch relations from 32 development source clusters.
On the same raw 64-candidate support, exact relation coverage at rank one
increases from 1.25\% for raw cosine to 7.73\% for the contextual D1 descriptor
and 8.56\% for its scalar score. Re-ranking cannot change coverage at rank 64
on this fixed support. Member-posterior proposals increase exact coverage at
rank 64 from 50.76\% to 73.36\%, with a paired exploratory interval of
[15.99,30.12] percentage points for the change. The matching checkpoint is
identical across these diagnostics and the current conditioning candidate;
production parity was checked separately.

For the member proposals, 31,393 relations have a wrong Top-1 but a GT match
inside the retained Top-64: 62.17\% of all observable relations and 70.00\% of
those missed by Top-1. This motivates retaining candidate uncertainty. It does
not prove that a learned attention head uses the candidates effectively or
improves generated assets. The denominator excludes unobservable relations,
and the measurement unit is a source ray paired with a target patch, not a
unique 3D point or a pixel. These historical matching diagnostics and the
SS128/cascade32 experiments have different measurement scopes and are not
treated as one paired downstream experiment.

The fixed-evidence candidate improves development SS macro IoU from 0.467355
for vanilla fine-tuning to 0.473471. Anchored learned probability refinement
reaches 0.474794 in the completed SS128 three-noise follow-up, a change of
0.001323 from fixed evidence. The source-cluster interval adjusted for the
two module contrasts is [-0.000409,0.003161], so this positive point estimate
does not establish the necessity of refinement. Its change from vanilla is
0.007439 with the corresponding two-contrast pipeline interval
[0.002585,0.012504]. These are development results from one training seed;
the learned candidate's actual cascade and independent confirmation are pending.

In the actual DEV32 fixed-evidence composition, however,
its geometry and appearance move in different directions. Relative to the
retargeting composition, Chamfer distance decreases by 0.00008579 and surface
F-Score increases by 0.005146, while source LPIPS increases by 0.001186 and
DreamSim increases by 0.005393. The exploratory paired interval for the LPIPS
change is [0.000224, 0.002275]; the F-Score interval is
[-0.000163, 0.011128]. Therefore this candidate is not established as a general
improvement over retargeting. It also illustrates why stage-level occupancy
accuracy is insufficient for selecting a complete editing pipeline.

The contextual matching descriptor has a positive SS point estimate relative
to raw semantic descriptors, whereas adding the scalar matching score has a
negative point estimate in that tested Top-1 path. Fixed Top-K aggregation and
additional learned Q/K refinement are consequently separate hypotheses.
Earlier learned attention improves source CLIP in one exploratory rendering
comparison while LPIPS/DreamSim do not improve. These results support reporting
endpoint-specific tradeoffs rather than a universal preservation claim.

The PBR same-geometry comparison does not currently establish an additional
benefit from feature correction or raw Top-K appearance input over Top-1
retargeting. Stage-specific composition and independent training confirmation
remain required before selecting the final pipeline.

\section{Ablation and final-table completion}
\begin{table}[t]
\centering\scriptsize
\resizebox{\linewidth}{!}{%
\begin{tabular}{lrrrrrrrrrr}
\toprule
&\multicolumn{3}{c}{Source preservation}&\multicolumn{2}{c}{Global naturalness}&\multicolumn{5}{c}{Target alignment}\\
Model&LPIPS$\downarrow$&DreamSim$\downarrow$&CLIP$\uparrow$&FID$\downarrow$&KID$\downarrow$&PSNR$\uparrow$&SSIM$\uparrow$&M-IoU$\uparrow$&CD$\downarrow$&F-Score$\uparrow$\\
\midrule
Zero-shot, source-only four views&\pending&\pending&\pending&\pending&\pending&\pending&\pending&\pending&\pending&\pending\\
Zero-shot, edited target + three sources&\pending&\pending&\pending&\pending&\pending&\pending&\pending&\pending&\pending&\pending\\
Vanilla V1&\pending&\pending&\pending&\pending&\pending&\pending&\pending&\pending&\pending&\pending\\
Retargeting&\pending&\pending&\pending&\pending&\pending&\pending&\pending&\pending&\pending&\pending\\
Selected combined model (if distinct)&\pending&\pending&\pending&\pending&\pending&\pending&\pending&\pending&\pending&\pending\\
\bottomrule
\end{tabular}}
\caption{Final TEST508 template. Every entry is pending. Remove a duplicate
combined-model row if selection produces an existing composition. Development
measurements must not be copied into this table.}
\label{tab:final-main-pending}
\end{table}

\begin{table}[t]
\centering\small
\begin{tabular}{lrrrrr}
\toprule
Matched final-context model&CD$\downarrow$&F-Score$\uparrow$&LPIPS$\downarrow$&DreamSim$\downarrow$&SS IoU$\uparrow$\\
\midrule
Selected full composition&\pending&\pending&\pending&\pending&\pending\\
Original spatial correspondence&\pending&\pending&\pending&\pending&\pending\\
Raw semantic descriptor&\pending&\pending&\pending&\pending&\pending\\
Actual Top-1 correspondence only&\pending&\pending&\pending&\pending&\pending\\
Individual voxel query&\pending&\pending&\pending&\pending&\pending\\
Without retargeting module&\pending&\pending&\pending&\pending&\pending\\
Mean descriptor (if member design retained)&\pending&\pending&\pending&\pending&\pending\\
Fixed refinement (if learned head retained)&\pending&\pending&\pending&\pending&\pending\\
\bottomrule
\end{tabular}
\caption{Final-context ablation template; all entries pending. Reuse an existing
arm only after confirming the same architecture, initialization, budget,
inputs, and evaluation. Conditional rows are omitted if absent from the
selected design. An optional same-Top-K score control is required only for
an isolated scalar-score claim.}
\label{tab:final-ablation-pending}
\end{table}

Only a missing matched control in the selected architecture requires new
ablation training. An inference-time module deletion is a diagnostic, not a
replacement for matched training. Top-1 versus weighted Top-K identifies a
combined candidate-count/aggregation effect; isolating scalar score utility
requires the same Top-K support and value path. Likewise, differences between
the original FULL implementation and a shared-evidence implementation cannot
be attributed solely to shared routing without a matching routing control.

\section{Limitations and pending claims}
Final ray-versus-voxel and mean-versus-member results, cascade confirmation
of learned versus fixed evidence, the final-context matched controls, independent training seed
confirmation, and held-out test results are pending. Geometric improvement
does not ensure source appearance improvement. Preserving the direct target
stream does not guarantee unchanged target output. Input-alpha support is not
surface visibility, and the source-use gate is not a calibrated occlusion
predictor. The final abstract, concluding superiority claims, quantitative
contribution statements, and final model label must follow these experiments.
\bibliographystyle{plain}
\bibliography{references}
\end{document}
